\chapter{硅\texorpdfstring{\quad}{ }胶体}

\section{碳族元素}
碳族元素属于元素周期表的第 \MyRoman{5} 主族，包括碳、硅、锗、 锡、铅五种元素。
碳族元素的原子序数、原子量和它们的原子的电子层排布如\cref{tab:4-1} 所示。
\begin{sidewaystable}
  \caption{碳族元素的原子序数、原子量和原子的电子层排布}\label{tab:4-1}
  \begin{tblr}{colspec={*4{c}lllX[l]ll},row{1,2}={m,c},hline{3}=0.8pt}
    \SetCell[r=2]{m,c} {元素名称} & \SetCell[r=2]{m,c} {元素符号} & \SetCell[r=2]{m,c} {原子序数} & \SetCell[r=2]{m,c} {原子量} & \SetCell[c=6]{m,c} 原子核外电子层排布 &&&&& \\
    & & & & K  & L  & M & N & O & P\\
  碳 & \ce{C}  & 6  & 12.01 & $1s^2$ & $2s^22p^2$ &  &  &  &  \\
  硅 & \ce{Si} & 14 & 28.09 & $1s^2$ & $2s^22p^6$ & $3s^23p^2$ &  &  &  \\
  锗 & \ce{Ge} & 32 & 72.59 & $1s^2$ & $2s^22p^6$ & $3s^23p^63d^{10}$ & $4s^24P^2$ &  &  \\
  锡 & \ce{Sn} & 50 & 118.7 & $1s^2$ & $2s^22p^6$ & $3s^23p^63d^{10}$ & $4s^24P^64d^{10}$ & $5s^25p^2$ &  \\
  铅 & \ce{Pb} & 82 & 207.2 & $1s^2$ & $2s^22p^6$ & $3s^23p^63d^{10}$ & $4s^24P^64d^{10}4f^{14}$ & $5s^25p^65d^{10}$ & $6s^26p^2$ \\
  \end{tblr}

  \smallskip
  \caption{碳族元素的一些重要性质}\label{tab:4-2}
  \begin{tblr}{colspec={ccccX[l]ccc},hline{3}=0.8pt,row{1,2}={m,c}}
    \SetCell[r=2]{m,c} {元\\素} & \SetCell[r=2]{m,c} {原子半径\\（\qty{e-10}{m}）} & \SetCell[r=2]{m,c} {第一电离能\\（\unit{eV}）} & \SetCell[r=2]{m,c} {主要\\化合价} & \SetCell[c=4]{m,c} 单质的性质 & & & \\
     & & &  & 状态和颜色 & {密度（\unit{g/cm^3}）}& {熔点（\unit{\celsius}）}& {沸点（\unit{\celsius}）}\\
    碳 & 0.77 & 11.3 & {$+2$，$+4$} & 无色或黑色固体   & {3.51\footnotemark[1] \\ 2.25\footnotemark[2] } & {3550 \\ \numrange{3652}{3697}（升华）} & {4827 \\ 4827} \\
    硅 & 1.17 &  8.2 & {$+4$}       & 从无色到棕色固体 & \numrange{2.32}{2.34} & 1410  & 2355\\
    锗 & 1.22 &  7.9 & {$+2$，$+4$} & 灰白色固体       &  5.35                 & 937.4 & 2330\\
    锡 & 1.41 &  7.3 & {$+2$，$+4$} & 银白色固体       &  7.28                 & 231.9 & 2260 \\
    铅 & 1.75 &  7.4 & {$+2$，$+4$} & 蓝白色固体       & 11.34                 & 327.5 & 1740 \\
  \end{tblr}
  \begin{minipage}{\linewidth}
    \footnotesize \cnum{1} 金刚石 \cnum{2} 石墨
  \end{minipage}
\end{sidewaystable}

第 \MyRoman{5} 主族元素位于周期表里容易失去电子的主族元素和容易得到电子的主族元素的中间位置，容易生成共价化合物。

碳族元素原子的最外层电子排布式是 $ns^2np^2$，随着电子层的增加，碳族元素的性质呈规律性的变化。\cref{tab:4-2} 列出了碳族元素的一些重要性质。

碳族元素随着电子层和核电荷数的增加，它们的一些重要物理性质和化学性质都发生规律性的变化。
它们从上到下由非金属性向金属性递变的趋势比氮族元素更为明显。
碳是明显的非金属；硅虽然外貌象金属，但在化学反应中更多地显非金属性，通常被认为是非金属；锗的金属性比非金属性强；锡和铅都是金属。

碳族元素的电负性比同一周期的卤族、氧族和氮族的小。

碳族元素的化合价除 $+4$ 价外，还有 $+2$ 价。
碳、硅、锗、锡的 $+4$ 价化合物是稳定的，而铅的 $+2$ 价化合物是稳定的。

在自然界里，碳族元素以不同的形态存在。
碳有游离态的碳，如金刚石、石墨、无定形碳等，也有化合态的碳，如各种碳酸盐和大量的有机化合物中的碳。
地壳中碳的含量虽然不多，但它是地球上形成化合物种类最多的元素。
硅在地壳中主要以含氧化合物矿石的形式存在。
锗是不常见的元素，通常它与若干金属同时存在于硫化物矿内。
锡和铅两种金属在自然界都以化合态存在。
锡的主要矿石是锡石矿（\ce{SnO2}），铅的主要矿石是方铅矿（\ce{PbS}）。

在初中化学里我们已学习过碳元素，下面将着重介绍硅及其重要化合物，锗、锡、铅就不再介绍了。

\begin{Practice}[习题]
\begin{question}
  \item 画出硅的原子结构简图，写出硅的电子式以及主要化合价，说明硅在自然界的存在方式。
  \item 说明碳、硅的化合物为什么大多数以共价键相结合。
\end{question}
\end{Practice}

\section{硅及其重要的化合物}
\subsection{硅}
我们已经知道，在地壳里，硅的含量在所有元素中居第二位，仅次于氧。
自然界里没有游离态的硅，而只有化合态的硅。
化合态的硅几乎全部是二氧化硅和硅酸盐，它们广泛地存在于地壳的各种矿物和岩石里。
硅是构成矿物与岩石的主要元素。

\subsubsection{物理性质}
晶体硅是无色到棕色有金属光泽、硬而脆的固体（\cref{tab:4-2}）。
硅原子最外电子层的结构是 $3s^23p^2$，有 4 个价电子。
晶体硅的结构与金刚石结构相似，每个硅原子跟另外 4 个硅原子形成 4 个共价键，成为正四面体结构。
这种结构跟金刚石结构相似，也是一种网状的原子晶体（\cref{fig:4-1}）\footnote{\cref{fig:4-1} 是平面图，晶体硅实际上是立体结构。}。
这样就决定了硅的硬度较大，熔点和沸
\par\medskip\noindent
\begin{minipage}{0.6\linewidth}\parindent2em
\noindent 点较高。
但是由于硅的原子半径比碳大，晶体硅里 \ce{Si-Si} 的键长（\qty{2.35e-10}{m}）大于金刚石里 \ce{C-C} 的键长；而 \ce{Si-Si} 的键能（\qty{42.5}{kCal/mol}）小于金刚石里 \ce{C-C} 的键能。
所以，晶体硅的硬度、熔点和沸点都小于或低于金刚石。

硅的导电性能介于金属和绝缘体之间。
硅是良好的半导体材料，可用来制造半导体器件，如硅整流器、晶体管和集成电路等。

锗和硅相似，也是重要的半导体材料。
\end{minipage}\hfill
\begin{minipage}{0.38\linewidth}
\begin{figurehere}
  \includegraphics{4-1.pdf}
  \caption{晶体硅结构的平面示意图}\label{fig:4-1}
\end{figurehere}
\end{minipage}



\subsubsection{化学性质}
硅是非金属元素，硅的许多化学性质和碳相似，它跟其它元素化合时形成共价键。

硅的化学性质不活泼。在常温下，除氟气、氢氟酸和强碱溶液外，其它物质如氧气、氯气、硫酸和硝酸等都不跟硅起反应。
在加热条件下，硅能跟一些非金属起反应，例如，把硅研细后加热，它就燃烧生成二氧化硅，同时放出大量的热。
\[ \ce{ Si + O2 \xlongequal{\triangle} SiO2 }\]

硅能跟强碱溶液作用生成硅酸盐，放出氢气。
\[ \ce{Si + 2NaOH + H2O \xlongequal{\quad} Na2SiO3 + 2H2 ^} \]

在一般情况下，硅跟氢气不能直接化合。硅的氢化物常用间接方法制得。

象碳一样，硅还能跟某些金属生成硅化物，所以，硅可用来制造合金。
如炼钢时常用硅铁合金做脱氧剂来除氧；含硅 4\% 的钢有导磁性，可以用来制造变压器的铁芯；含硅 15\% 左右的钢有耐酸性，可以用来制造耐酸设备。

工业上，硅是在电炉里用碳还原二氧化硅而制得的。
\[ \ce{ SiO2 + 2C \xlongequal{\text{高温}} Si + 2CO ^}\]

这样制得的硅是含少量杂质的粗硅。
粗硅提纯后，可制得作为半导体材料的高纯硅。

\subsection{二氧化硅}
\subsubsection{物理性质}
二氧化硅 \ce{SiO2} 是一种坚硬难熔的固体，它和其它矿物构成了多种岩石，广泛地分布在自然界里。

天然的二氧化硅分为晶体和无定形两大类。
石英的主要成分就是二氧化硅晶体。
自然界纯净的石英是无色透明的六方柱状晶体，这种晶体就是我们常说的水晶。

硅藻土含有无定形二氧化硅，它是死去的硅藻\footnote{硅藻是单细胞的低等水生植物。}及其它微小生物的遗体经沉积胶结而成的多孔、质轻、松软的固体物质。
它的表面积很大，吸附能力较强，可作吸附剂和催化剂的载体\footnote{为了增加催化剂的有效面积，一般使催化剂附着于多孔的物体表面，这种多孔物体叫做载体。}，以及保温材料等。

二氧化硅和二氧化碳在物理性质上有很大的差别。
例如，二氧化硅的熔点高、硬度大；而二氧化碳在通常状况下是气体，固体二氧化碳的熔点很低，等等。
二氧化硅和二氧化碳的这些差别是由它们的不同结构决定的。
我们已经知道固体二氧化碳是一种分子晶体。
分子间只有较小的分子间引力，所以，它的熔点很低。
但是二氧化硅却不是由单个的 “\ce{SiO2}” 的分子所组成的分子晶体，而是一种原子晶体。
如\cref{fig:4-2}\footnote{\cref{fig:4-2} 是一个简化的平面示意图，实际上二氧化硅晶体是立体的网状结构。}所示，1 个 \ce{Si} 原子跟 4 个 \ce{O} 原子形成了 4 个共价键，这样，每 1 个 \ce{Si} 原子周围结合 4 个 \ce{O} 原子；同时，每个 \ce{O} 原子跟两个 \ce{Si} 原子相结合。
实际上，二
\par\medskip\noindent
\begin{minipage}{0.55\linewidth}\parindent2em
氧化硅晶体是由硅原子和氧原子按 $1:2$ 的比率所组成的立体网状的原子晶体，我们通常用化学式 \ce{SiO2} 来表示二氧化硅的组成。
在二氧化硅的晶体里 \ce{Si-O} 的键长为 \qty{1.62e-10}{m}，键能为 \qty{88.2}{kCal/mol}。
由于二氧化硅晶体里 \ce{Si-O} 键的键能很高，并形成了一种立体网状的原子晶体，所以，要使它熔融，也就是说要破坏二氧化硅的晶体，必须消耗较多的能量，因此，二氧化硅的熔点很高，硬度也很大。
\end{minipage}\hfill
\begin{minipage}{0.42\linewidth}
\begin{figurehere}
  \includegraphics{4-2.pdf}
  \caption{二氧化硅晶体平面示意图}\label{fig:4-2}
\end{figurehere}
\end{minipage}

\subsubsection{化学性质}
由于二氧化硅的 \ce{Si-O} 键能很大，因而它的化学性质十分稳定，不能跟酸（除氢氟酸外）发生反应。
二氧化硅是一种酸性氧化物。
但是，二氧化硅不溶于水，它不能跟水起反应生成酸。
二氧化硅能跟碱性氧化物或强碱反应生成盐。
\begin{gather*}
  \ce{ SiO2 + CaO \xlongequal{\text{高温}} CaSiO3 }\\
  \ce{ SiO2 + 2NaOH \xlongequal{\quad} Na2SiO3 + H2O }
\end{gather*}

玻璃中含有二氧化硅，因而它能被碱液腐蚀。
实验室中盛放碱液的试剂瓶常用橡皮塞，而不用玻璃塞，就是为了防止玻璃受碱液腐蚀生成 \ce{Na2SiO3}，而使瓶口和塞子粘在一起。

\subsubsection{用途}
二氧化硅的用途很广。
自然界里比较稀少的水晶可用以制造电子工业的重要部件、光学仪器和工艺品。

二氧化硅是制造光导纤维的重要原料。

一般较纯净的石英，可用来制造石英玻璃。
石英玻璃的膨胀系数很小,相当于普通玻璃的 1/18，能经受温度的剧变，耐酸性能好（除 \ce{HF} 外）。
因此，石英玻璃常用来制造耐高温的化学仪器。

石英砂常用作玻璃原料和建筑材料。


\subsection{硅酸\texorpdfstring{\quad}{ }硅酸盐}
\subsubsection{硅酸}
硅酸不能用二氧化硅跟水直接作用制得，而只能用相应的可溶性的硅酸盐跟酸作用制得。

硅酸钠（\ce{Na2SiO3}）\footnote{\ce{Na2SiO3} 通常也叫偏硅酸钠，现在根据 1982 年版《无机化学命名原则》称为硅酸钠。}的水溶液跟盐酸起反应生成的白色胶状沉淀叫做原硅酸，通常用 \ce{H4SiO4} 来表示它的组成。
原硅酸几乎不溶于水，是一种弱酸，很不稳定。
这种白色胶状物在空气里干燥，失去一部分水后，变成白色粉末。
这种物质是硅酸，通常用 \ce{H2SiO3} 来表示它的组成（可以认为是 \ce{H4SiO4 = H2SiO3 + H2O}。硅酸不溶于水，也是一种弱酸，它的酸性比碳酸还弱。

\subsubsection{硅酸盐}
硅酸、原硅酸以及由它们缩水结合而成的各种酸所对应的盐，统称硅酸盐。
例如，硅酸钠（\ce{Na2SiO3}）、镁橄榄石（\ce{Mg2SiO4}）、高岭石（\ce{Al2(Si2O5)(OH)4}）等都是硅酸盐。
硅酸盐种类很多，结构也很复杂，它是构成地壳岩石的最主要成分。
我们可以用二氧化硅和金属氧化物的形式表示硅酸盐的组成。
例如：

\begin{tblr}{colspec={ll},hlines=0pt,vlines=0pt}
  硅酸钠   & \ce{Na2SiO3} \quad （\ce{Na2O.SiO2} ）\\
  镁橄榄石 & \ce{Mg2SiO4} \quad （\ce{2MgO.SiO2}）\\
  高岭石   & \ce{Al2(Si2O5)(OH)4} \quad （\ce{Al2O3.2SiO2.2H2O}）\\
\end{tblr}

许多硅酸盐都是难溶于水的。
可溶性硅酸盐中，最常见的是 \ce{Na2SiO3}，它的水溶液俗名水玻璃。
水玻璃是无色粘稠的液体，是一种矿物胶，它既不能燃烧也不受腐蚀，在建筑工业上可用作粘合剂、耐酸水泥掺料等。
木材、织物浸过水玻璃后，具有防腐性能，且不易着火。
水玻璃还可用作耐火材料。

粘土的成分也是硅酸盐。
花岗岩里的正长石（\ce{KAlSi3O8}）在二氧化碳和水的作用下，分解而生成粘土等物质。
粘土的种类很多，成分也很复杂（主要是高岭石），它是土壤里矿物质的主要部分。
常见的有高岭土\footnote{高岭土又叫瓷土，主要由高岭石的微细晶体组成，因盛产于我国江西景德镇的高岭而得名。}和一般粘土，前者含杂质较少，后者含杂质较多。

\begin{Practice}[习题]
\begin{question}
  \item 粗硅提纯时常用下述方法: 将粗硅在高温下跟氯气反应生成四氯化硅 \ce{SiCl4}，\ce{SiCl4} 经过分馏提纯，再用氢气还原得到纯硅。写出上述反应的化学方程式。
  \item 晶体硅和金刚石在结构和物理性质上有哪些相似的地方，有哪些不同的地方？
  \item 二氧化碳和二氧化硅在物理性质方面有什么不同，试从结构角度加以说明。
  \item 怎样用化学方法检验生石灰里混有的石英和石灰石等杂质？写出有关的化学方程式。
  \item 往电炉里加入 \qty{60}{g} \ce{SiO2} 和 \qty{20}{g} 碳的混和物，通电，使它们发生如下反应: $\ce{SiO2 + 2C \xlongequal{\quad} Si + 2CO}$。求：
  \begin{tasks}
    \task 生成物的质量各多少克？
    \task 生成的 \ce{CO} 在标准状况下是多少升？
  \end{tasks}
  \item 把下列式子改写成氧化物的形式：
  \begin{tasks}
    \task 高岭石 （\ce{Al2(Si2O5)(OH)4}）
    \task 滑石   （\ce{Mg3(Si4O10)(OH)2}）
    \task 钙沸石 （\ce{Ca(Al2SI3O10).3H2O}）
  \end{tasks}
  \item 怎样用石英和其它物质制取硅酸？ 写出相应的化学方程式。
\end{question}
\end{Practice}



\section{硅酸盐工业简述}
以含硅物质为原料，经过加热制成硅酸盐产品的工业，如制造水泥、玻璃、陶瓷等产品的工业，叫做硅酸盐工业。
它在国民经济中占有重要的地位。

\subsection{水泥}
普通硅酸盐水泥的主要原料是石灰石和粘土。
先把石灰石、粘土和其它辅助原料按一定的比率混和，磨细成生料，将生料装入回转窑里煅烧（\cref{fig:4-3}）。
\begin{figure}
  \includegraphics{4-3.pdf}
  \caption{水泥回转窑示意图}\label{fig:4-3}
\end{figure}

煅烧时，燃料从回转窑的低端喷入，原料从回转窑的高端流入。
借原料本身的重力作用逐渐下移，经过预热带、煅烧带和烧结带（温度可达 \qtyrange{1400}{1500}{\celsius}）。
原料在高温下，发生了复杂的物理、化学变化，成为部分熔化状态，冷却后成为硬块，这种物质叫做熟料，再加入适量的石膏（调节水泥硬化速度），磨成细粉，即制成普通硅酸盐水泥。
这种水泥的主要成分有：

\begin{tblr}{colspec={ll},hlines=0pt,vlines=0pt}
  硅酸三钙 & \ce{3CaO.SiO2}\\
  硅酸二钙 & \ce{2CaO.SiO2}\\
  铝酸三钙 & \ce{3CaO.Al2O3}\\
\end{tblr}

水泥实际上是上述主要成分的混和物，水泥的组成和结晶形态的不同直接影响到它的各种主要性能。

水泥具有水硬性。
水泥跟水拌和后，发生作用，生成不同的水合物，同时放出一定的热量。
生成的水合物逐步形成胶状物，并开始凝聚，最后，有些胶状物转变为晶体，使胶状物和晶体交错地结合起来，成为强度很大的固体。
这个过程叫做水泥的硬化。
水泥不论是在空气中还是在水中都能硬化，所以，它不仅是一般的建筑材料，而且还是水下工程必不可少的建筑材料。

为了改善水泥的性能，扩大水泥的使用范围，可在硅酸盐水泥熟料里，掺入适当比率的混和材料，制成各种水泥。
例如，矿渣硅酸盐水泥就是在硅酸盐水泥熟料里加入一定量的高炉矿渣（主要成分是 \ce{CaSiO3}），沸石岩水泥就是在水泥熟料里掺入一定量的沸石岩。

水泥、砂子和水的混和物叫做水泥砂浆，在建筑上用它作粘合剂，能把砖、石等粘结起来。
水泥、砂子、碎石和水按一定比率的混和物硬化后叫做\Concept{混凝土}，常用它来建筑桥梁、厂房等巨大建筑物。
水泥的热膨胀系数几乎跟铁一样，所以用混凝土建造建筑物常用钢筋作结构，使建筑物更加坚固，这叫做钢筋混凝土。

\subsection{玻璃}
制造普通玻璃的主要原料是纯碱（\ce{Na2CO3}）、石灰石（\ce{CaCO3}）和石英（\ce{SiO2}），有些特种玻璃原料中还包含氧化铅（\ce{PbO}）和硼砂（\ce{Na2B4O7.10H2O}）。
生产玻璃时，把原料粉碎，按适当比混和以后，放入玻璃熔炉里，加强热。
原料熔融后发生了比较复杂的物理、化学变化，其中主要的反应是二氧化硅跟碳酸钠和碳酸钙起反应生成硅酸盐和二氧化碳：
\begin{gather*}
  \ce{ Na2CO3 + SiO2 \xlongequal{\text{高温}} Na2SiO3 + CO2 ^} \\
  \ce{ CaCO3 + SiO2 \xlongequal{\text{高温}} CaSiO3 + CO2 ^} 
\end{gather*}

在原料里，石英的用量是较多的。
所以，普通玻璃是 \ce{Na2SiO3}、\ce{CaSiO3} 和 \ce{SiO2} 熔化在一起所得到的物质。
这种物质不是晶体，称作玻璃态物质，它没有一定的熔点，而是在某一温度范围内逐渐软化。
在软化状态时，玻璃可以制成任何形状的制品。

玻璃的种类很多，除上述普通玻璃外，还有硼酸盐玻璃，它是耐热仪器玻璃，膨胀系数小，能够耐骤冷、骤热，化学稳定性好，能耐酸、碱的腐蚀，可用于制造化学仪器。
铅玻璃折光率强，常用于制作光学仪器等。
制造有色玻璃，一般是在原料里加入某些金属氧化物，使金属氧化物均匀地分散到玻璃态的物质里，使玻璃呈现出特征颜色。
例如，加氧化钴（\ce{Co2O3}）呈蓝色；加氧化亚铜（\ce{Cu2O}）呈红色，等等。

把普通玻璃放入电炉里加热，使它软化，然后急速冷却，得到钢化玻璃。
钢化玻璃的机械强度比普通玻璃大 \numrange{4}{6} 倍，不易破碎。
破碎时碎块没有尖锐的棱角，不易伤人，可用来制造汽车或火车的车窗等。

玻璃还可以制成纤维。
它具有较高的强度，可作隔音、隔热、电气绝缘材料，也可制造玻璃纤维增强塑料等等。

\begin{Practice}[习题]
\begin{question}
  \item 制造普通玻璃的主要原料在玻璃窑里熔化后，发生了哪些主要化学反应？写出反应的化学方程式。
  \item\label{prac:4-3-2} 普通玻璃含有 13\% 的 \ce{Na2O}、11.7\% 的 \ce{CaO}、75.3\% 的 \ce{SiO2}。试计算三种氧化物的物质的量之比。
  \item 制造一吨普通玻璃，需要纯碱、石灰石和二氧化硅各多少吨（参看第 \ref{prac:4-3-2} 题得到的物质的量之比）？并计算生产一吨普通玻璃，在温度为 \qty{1200}{\celsius}、压强为 \qty{750}{mmHg} 时，产生的二氧化碳是多少升？
\end{question}
\end{Practice}


\section{胶体}
在学习前几节时，我们已经涉及有些物质在一定条件下形成了胶体。
例如，硅酸在水里形成的透明液体；某些金属氧化物分散到玻璃态物质里形成有色玻璃等等。
本节主要介绍什么是胶体和胶体的重要性质。
\subsection{胶体}
在初中化学里我们已经学习过溶液、悬浊液和乳浊液。
它们都是一种物质（或几种物质）的微粒分散于另一种物质里形成的混和物。
我们常把这种混和物叫做分散系。
其中分散成微粒的物质叫做分散质；微粒分散在其中的物质叫做分散剂。
例如，对溶液来说，溶质是分散质，溶剂是分散剂，溶液是一种分散系。
悬浊液和乳浊液也是一种分散系，其中的固体小颗粒或小液滴是分散质，所用的溶剂是分散剂。

胶体也是一种分散系，在这种分散系里，分散质微粒直径的大小介于溶质的分子或离子的直径（一般小于 \qty{e-9}{m}）和悬浊液或乳浊液微粒的直径（一般大于 \qty{e-7}{m}）之间。
一般地说，\emph{分散质微粒的直径大小在 \qtyrange{e-10}{e-7}{m} 之间的分散系，叫做}\Concept{胶体}。

\begin{Experiment}
往一个烧杯里注入蒸馏水 \qty{20}{ml}，给烧杯加热使水沸腾，然后往沸水里滴加氯化铁（\ce{FeCl3}）饱和溶液 \qtyrange{1}{2}{ml}。
继续煮沸，待溶液呈红褐色后，停止加热。
观察制得的胶体。

取一个大试管，注入 $0.01M$ 碘化钾（\ce{KI}）的溶液 \qty{10}{ml}，用胶头滴管滴入 \numrange{8}{10} 滴相同浓度的硝酸银（\ce{AgNO3}）溶液，边滴入边振荡。观察制得的胶体。
\end{Experiment}

上面两个实验可分别用化学方程式表示如下：
\begin{gather*}
  \ce{ FeCl3 +3H2O \xlongequal{\triangle} Fe(OH)3 \text{（胶体）} + 3HCl }\\
  \ce{ KI +AgNO3 \xlongequal{\triangle} AgI \text{（胶体）} + KNO3 }
\end{gather*}

从这两个实验里可以分别看到，氢氧化铁分散到水里所形成的胶体是红褐色的；碘化银分散到水里所形成的胶体是浅黄色的。

这两种胶体的微粒分别是氢氧化铁或碘化银的许多分子的集合体\footnote{这种胶体也叫粒子胶体}，它们的直径大小在 \qtyrange{e-9}{e-7}{m} 之间。

有些物质的分子直径很大，达到了胶体微粒的大小，并且能溶解于水（或其他溶剂）。这种物质溶于水（或其他溶剂），就形成了胶体，这种胶体的微粒是分子\footnote{这种胶体也叫分子胶体}。

我们常常根据胶体粒子大小介于 \qtyrange{e-9}{e-7}{m} 之间这一特点，把不纯的胶体放进有半透膜\footnote{半透膜: 一般指动物的膀胱膜、肠衣、羊皮纸、胶棉薄膜、玻璃纸等。半透膜有非常小的细孔，这些细孔只能使离子或分子透过而不能使胶体微粒透过。}的容器内，让较小的微粒（如分子、离子）通过半透膜，以净化胶体。\pagebreak

\begin{Experiment}*[righthand ratio=0.4]
  用半透膜制成一个袋，往这个袋里注入淀粉胶体 \qty{10}{ml} 和食盐溶液 \qty{5}{ml} 的混和液体，然后如\cref{fig:4-4} 所示用线把半透膜袋的上口缚好，系在玻璃棒上，并把它悬挂在盛有蒸馏水的烧杯里。

  几分钟后，用两个试管各取烧杯里的液体 \qty{5}{ml}。往其中一个试管里注入少量硝酸银溶液；往另一个试管里注入少量碘水。观察这两个试管里所发生的变化。
  \tcblower 
  \begin{figurehere}
    \includegraphics{4-4.pdf}
    \caption{渗析}\label{fig:4-4}
  \end{figurehere}
\end{Experiment}

从实验可以看到，在加入硝酸银溶液的试管里出现了白色沉淀；在加入碘水的试管里不发生变化。
这就证明了氯离子可以透过半透膜的微孔，而淀粉胶体的微粒不能透过半透膜的微孔。
同样可以用实验证明，钠离子也能透过半透膜。
这种现象说明胶体微粒大于溶液里溶质的离子或分子。
象这种把混有离子或分子杂质的胶体装入半透膜的袋里，并把这个袋放在溶剂中，从而使离子或分子从胶体溶液里分离的操作叫做\Concept{渗析}。
应用渗析的方法可精制某些胶体。

胶体的种类很多，按照分散剂的不同，可分为液溶胶、气溶胶和固溶胶。
分散剂是液体的，叫做\Concept{液溶胶}（也叫\Concept{溶胶}），例如,上面实验里制备的 \ce{Fe(OH)3} 和 \ce{AgI} 胶体都是液溶胶。
分散剂是气体的，叫做\Concept{气溶胶}，例如，雾、云、烟等都是气溶胶。
分散剂是固体的，叫做\Concept{固溶胶}，例如，烟水晶、有色玻璃等都是固溶胶。

\subsection{胶体的重要性质}
从外观看胶体跟溶液并没有明显区别。
胶体的微粒虽然比溶液里溶质的分子或离子大很多，但是它们都能通过滤纸。 
那么，怎样来鉴别胶体和溶液呢？ 
下面介绍胶体的几种重要的性质，其中某些性质可用来鉴别胶体和溶液。
\subsubsection{丁达尔\protect\footnote{丁达尔（J. Tyndall 1820--1893），英国物理学家。}现象}
当太阳光透过窗户上的小孔射到屋里的时候，从入射光线垂直的方向观察，可以看到一条光亮的 “通路”，这种现象叫做光的散射。
因为光束在空气里前进时，遇到很多灰尘的微粒，它们的直径很小，使光束的部分光线偏离原来的方向而分散传播。
如果让光线透过胶体，从侧面同样可以观察到胶体里也出现一条光亮的“通路”。
这也是由于胶体微粒对光线的散射而形成的（\cref{fig:4-5}）。
这个现象叫做\Concept{丁达尔现象}。

\begin{figure}
  \begin{minipage}[b]{0.5\linewidth}\centering
    \includegraphics[height=4cm]{4-5.jpg}\par
    {\footnotesize 左侧烧杯里是溶液，右侧烧杯里是胶体。}
    \caption{丁达尔现象}\label{fig:4-5}
  \end{minipage}
  \begin{minipage}[b]{0.48\linewidth}\centering
    \includegraphics{4-6.pdf}\par
    {\footnotesize 图中每条折线是一个微粒在每隔一定时间所在位置的连线。}
    \caption{布朗运动示意图}\label{fig:4-6}
  \end{minipage}
\end{figure}

当光线通过溶液时看不到这种现象。
可以用这种方法鉴别胶体和溶液。

\subsubsection{布朗\protect\footnote{布朗（R. Brown 1773--1858），英国植物学家。}运动}
1827 年，布朗把花粉悬浮在水里，用显微镜观察，发现花粉的小颗粒作不停的、无秩序的运动，这种现象叫做\Concept{布朗运动}（\cref{fig:4-6}）。

用超显微镜\footnote{超显微镜是一种特殊的显微镜，从显微镜侧面进入强光，由于胶体的丁达尔现象可以观察到胶体微粒的存在和动态。}观察溶胶，可见胶体微粒也在进行布朗运动。
这是因为水分子（或分散剂的分子）从各方面撞击胶体微粒，而每一瞬间胶体微粒在不同方向受的力是不相同的，所以胶体微粒运动的方向每一瞬间都在改变，因而形成不停的、无秩序的运动。

\subsubsection{电泳现象}
\par\noindent
\begin{minipage}{0.55\linewidth}\parindent2em
在一个 U 形管里盛有红褐色 \ce{Fe(OH)3} 胶体，从 U 形管的两个管口各插入一个电极（\cref{fig:4-7}）。
通直流电后，发现阴极附近的颜色逐渐变深，阳极附近的颜色逐渐变浅。
这表明 \ce{Fe(OH)3} 胶体微粒带正电荷，在电场的影响下，向阴极移动。
象这样在外加电场的作用下，胶体的微粒在分散剂里向阴极（或阳极）作定向移动的现象，叫做\Concept{电泳}。
\end{minipage}\hfill
\begin{minipage}{0.4\linewidth}
\begin{figurehere}
  \includegraphics{4-7.pdf}
  \caption{电泳现象}\label{fig:4-7}
\end{figurehere}
\end{minipage}\par\medskip

电泳现象证明了胶体的微粒是带有电荷的。
那么，为什么胶体的微粒带有电荷呢？
一般地说来，是由于胶体的微粒具有较大的表面积，能吸附阳离子或阴离子，因而带有正电荷或负电荷。

当物质分散成胶体的微粒时，分散质的总表面积有很大的增加。
例如，每边长为 \qty{1}{cm} 的立方体，它的总表面积为 \qty{6}{cm^2}。
当把这个物体分散成胶体的微粒时，如每边长为 \qty{e-8}{m} 的小立方体，它的总表面积即成为 \qty{600}{m^2}，也就是比原来增加了 100 万倍。
由于胶体的微粒有很大的表面积，所以具有较强的吸附能力。
不同的胶体微粒吸附不同电荷的离子。
有些胶体的微粒吸附阳离子，有些胶体的微粒吸附阴离子。
一般说来，金属氢氧化物、金属氧化物的胶体微粒吸附阳离子，胶体微粒带正电荷；非金属氧化物、金属的硫化物的胶体微粒吸附阴离子，胶体微粒带负电荷。
因为胶体的微粒是带电的粒子，所以，在电场的作用下，发生了定向运动，产生了电泳现象。

\subsubsection{胶体的凝聚}
由于同一种胶体微粒带有相同的电荷，胶体的微粒相互排斥，在一般情况下，胶体的微粒不容易聚集，因而胶体是比较稳定的分散系，可以保存较长的时间。
但是，如果往某些胶体里加入少量的电解质，由于电解质电离生成的阳离子或阴离子，中和了胶体微粒所带电荷，使胶体的微粒聚集成较大的颗粒，形成了沉淀，从分散剂里析出，这个过程叫做\Concept{凝聚}。

\begin{Experiment}
  在一个试管里盛 \ce{Fe(OH)3} 胶体 \qty{5}{ml}，滴入 \ce{MgSO4} 溶液 \qtyrange{1}{2}{ml}，振荡。观察发生的现象。
\end{Experiment}

从上面实验可以看到，\ce{Fe(OH)3} 胶体里加入 \ce{MgSO4} 溶液后，胶体变成浑浊状态，说明胶体微粒发生了凝聚作用。
这是因为 \ce{MgSO4} 溶液中的 \ce{SO4^2-} 离子，跟 \ce{Fe(OH)3} 胶体的微粒所带电荷发生了电中和，使胶体的微粒聚集成为沉淀析出。
除加入电解质可使某些胶体凝聚外，给胶体加热，或把两种带有相反电荷的胶体混和，也可使胶体发生凝聚作用。

在一般情况下，胶体发生凝聚作用都生成沉淀。
但有些胶体凝聚后，胶体的微粒和分散剂凝聚在一起成为不流动的冻状物，这种物质是一种凝胶。
例如，我们日常食用的豆腐就是把盐卤（它的主要成分是 \ce{MgCl2.6HO}） 或石膏（\ce{CaSO4.2H2O}）溶液加入豆浆里，使豆浆里的蛋白质和水等物质一起凝聚而制成的一种凝胶。

往硅酸钠溶液里加入盐酸，生成硅酸的沉淀。
在减压的情况下把这种沉淀加热至 \qty{300}{\celsius}，硅酸失去水分，变成了一种网状多孔物质。
它的主要成分是约含 4\% 水的二氧化硅，这种物质叫做硅胶。
硅胶也是一种凝胶，它的表面积很大，吸附能力较强，可作干燥剂、吸附剂以及催化剂的载体。

胶体的知识对工农业生产、科学研究、国防以及日常生活等都有十分重要的意义。
这是因为动植物的许多生理现象必须用胶体的知识来科学地解释。
土壤里许多物质如粘土、腐殖质等常以胶体形式存在，所以，在土壤里发生的一些化学过程跟胶体的关系很密切。
国防工业上有些火药、炸药必须制成胶体，冶金工业上的选矿，石油原油的脱水，塑料、橡胶以及合成纤维等的制造过程都应用到胶体的知识。
日常生活里经常接触和应用的胶体，有食品中的牛奶、豆浆、粥，用品中的塑料、橡胶制品，建筑材料中的水泥，等等。
由于胶体的知识在许多方面有很重要的用途，它已经发展成为一门独立的学科。

\begin{Practice}[习题]
\begin{question}
  \item 制备 \ce{Fe(OH)3} 胶体时，有一部分 \ce{FeCl3} 溶解在里面，怎样把它除去？检验除去 \ce{FeCl3} 后的胶体中是否仍含有 \ce{Cl-}，简述其过程，写出反应的离子方程式。
  \item 怎样用实验的方法鉴别溶液和胶体？
  \item 下列说明中哪一项能够正确地说明胶体粒子的电泳现象。
  \begin{tasks}
    \task 电泳现象是因胶体微粒带有正或负的电荷而引起的。
    \task 电泳现象的原因是液溶胶中一半粒子带正电，一半粒子带负电。
    \task 电泳现象是布朗运动的结果。
  \end{tasks}
  \item 在一定时间内，胶体为什么能够稳定存在？
  \item 硫化砷的胶体粒子因电泳而向阳极运动。在硫化砷胶体中加入下列物质是否会出现凝聚？
  \begin{tasks}(2)
    \task \ce{Fe(OH)3} 胶体
    \task \ce{AlCl3} 溶液
  \end{tasks}
  \item 为什么含有泥砂胶体颗粒的河水，跟海水相遇，容易沉积而形成沙洲（三角洲）？
  \item 列表比较溶液、胶体、悬浊液或乳浊液的不同点。
\end{question}
\end{Practice}

\section*{内容提要}
\addcontentsline{toc}{section}{内容提要}\setcounter{subsection}{0}
\subsection*{一、碳族元素}
碳族元素属于元素周期表的第 \MyRoman{5} 主族，包括碳、硅、锗、 锡、铅五种元素。碳族元素的最外层电子排布式是 $ns^2np^2$，碳族元素的电负性比同一周期的卤族、氧族和氮族元素小。
碳族元素位于元素周期表里容易失去电子的主族元素和容易得到电子的主族元素的中间位置。
碳族元素除 $+4$ 价外，还有 $+2$ 价。

碳有游离态的碳和化合态的碳两种形式，它是地球上形成化合物种类最多的元素。
硅在地壳中的含量居第二位，它主要以含氧化合物矿石形式存在。
锗常与若干金属同时存在于硫化物矿内，锡和铅都以化合状态存在于自然界。

\subsection*{二、硅及其重要化合物}\setcounter{subsubsection}{0}
\subsubsection{硅}
硅的原子序数为 14，电子排布式是 $1s^22s^22p^63s^23p^2$。
硅原子有 4 个价电子，能够形成 4 个共价键。

晶体硅的结构跟金刚石的结构相似，也是一种原子晶体。
晶体硅和锗是良好的半导体。

硅的化学性质不活泼。
在常温下，除氟气、氢氟酸和强碱溶液外，其它物质如氧气、氯气、硫酸和硝酸等都不跟硅起反应。
在加热条件下，硅跟氧气起反应。
硅还能跟强碱溶液作用，生成硅酸盐。

\subsubsection{二氧化硅}
天然的二氧化硅分为晶体和无定形两大类。
二氧化硅的晶体是由硅原子和氧原子组成的原子晶体，可用化学式 \ce{SiO2} 来表示它的组成。

二氧化硅不溶于水，是一种酸性氧化物，能跟碱性氧化物和碱起反应生成盐。

\subsubsection{硅酸和硅酸盐}
硅酸有原硅酸（\ce{H4SiO4}）、硅酸（\ce{H2SiO3}）等。
它们是一类弱酸，它们的酸性比碳酸还要弱。
可用硅酸钠跟盐酸起反应制得硅酸。

原硅酸、硅酸以及由它们缩水结合而成的酸所对应的盐统称硅酸盐。
如硅酸钠（\ce{Na2SiO3}）、高岭石（\ce{Al2(Si2O5)(OH)4}）等。

可溶性硅酸盐中，最常见的是 \ce{Na2SiO3}，其水溶液俗名水玻璃。

\subsubsection{硅酸盐工业}
用石灰石和粘土作主要原料可制造水泥。
用纯碱、石灰石和石英作原料可制造普通玻璃。

\subsection*{三、胶体}\setcounter{subsubsection}{0}
\subsubsection{胶体}
分散质微粒的直径在 \qtyrange{e-9}{e-7}{m} 之间的分散系，叫做胶体。
胶体的微粒不能透过半透膜，因此可用渗析方法净化胶体。
根据胶体分散剂的状态不同可把胶体分为气溶胶、液溶胶和固溶胶。

\subsubsection{胶体的性质}
胶体的性质有丁达尔现象、布朗运动和电泳现象等。
胶体的微粒在一定条件下（如加入少量电解质或加热等），可以聚集而发生沉淀，这个过程叫胶体的凝聚。

\begin{Review}
\begin{question}
  \item 硅和二氧化硅都属于哪一类晶体？它们在结构上和性质上有哪些相似点？
  \item 晶体硅的键长\CJKunderline[hidden]{\hspace{1em}大于\hspace{1em}}金刚石的键长，因而金刚石的键能\CJKunderline[hidden]{\hspace{1em}大于\hspace{1em}}晶体硅的键能，晶体硅的熔点比金刚石的熔点\CJKunderline[hidden]{\hspace{1em}低\hspace{1em}}。
  \item 比较 \ce{CO2} 和 \ce{SiO2} 在晶体结构上和物理性质上的主要不同点。
  \item 完成下列化学方程式，并注明反应条件。
  \begin{figurehere}
    \includegraphics{R4-4.pdf}
  \end{figurehere}
  \item 简要回答下列问题：
  \begin{tasks}
    \task 为什么水玻璃必须保存在密闭的容器里？
    \task 为什么盛氢氧化钠的试剂瓶，不能用玻璃塞？
  \end{tasks}
  \item 胶体溶液区别于其它分散系的本质特征是（\qquad）。
  \begin{tasks}[label=\Alph*.](2)
    \task 胶体微粒带电；
    \task 产生丁达尔现象；
    \task! 分散质微粒直径在 \qtyrange{e-9}{e-7}{m} 之间；
    \task 胶体微粒做布朗运动；
    \task 胶体微粒不能穿过半透膜。
  \end{tasks}
  \item 用氢氧化铁胶体作电泳实验，阳极周围颜色\CJKunderline[hidden]{\hspace{1em}变浅\hspace{1em}}。欲使氢氧化铁胶体凝聚可加入\CJKunderline*[hidden]{\hspace{1em}(1)、(3)、(4)\hspace{1em}}。
  \begin{tasks}(3)
    \task 硅酸胶体；
    \task \ce{Al(OH)3} 胶体；
    \task 硫化砷胶体；
    \task 硫酸镁溶液；
    \task 蒸馏水。
  \end{tasks}
\end{question}
\end{Review}